\documentclass[11pt,a4paper]{article}
\usepackage[margin=27mm]{geometry}
\usepackage{amsmath,amssymb,amsthm,hyperref}
\newtheorem{theorem}{Theorem}
\newtheorem{lemma}[theorem]{Lemma}
\newcommand{\F}{\mathbb F}
\newcommand{\Sz}{\operatorname{Sz}}
\title{A missing boundary in the Suzuki subfield construction}
\author{Human Compression project}\date{}
\begin{document}\maketitle
\begin{abstract}
The number of distinct prime divisors of a field degree does not determine
whether its smallest selected subfield is trivial. A repeated prime factor
allows all distinct primes to be selected while leaving a simple Suzuki
group at the bottom. This yields a rank-four residually weakly primitive
pregeometry omitted from the eleven literal constructions in the stated
classification. We give the complete subgroup argument, reusing the
earlier project proof of this boundary. We do not assert a corrected
exhaustive classification or a new solution of an open problem.
\end{abstract}

\section{The subgroup condition}
Let $G$ be a group, let $I$ be a finite type set, and let $G_i\le G$
for $i\in I$. The coset pregeometry has elements $gG_i$, of type $i$;
two elements are incident when their cosets intersect. Put
\[
 G_J=\bigcap_{j\in J}G_j\quad(J\ne\varnothing),\qquad
 G_\varnothing=G.
\]
The source calls this system residually weakly primitive when, for every
$J\subsetneq I$, at least one $i\in I\setminus J$ makes
$G_{J\cup\{i\}}$ a maximal proper subgroup of $G_J$.
This is the definition on printed page 4 of the source. Its principal
classification concerns pregeometries with precisely this subgroup
condition; flag transitivity and the geometry axiom are not additional
hypotheses of that classification.

Write $S_d=\Sz(2^d)$ for odd $d$. We use standard simultaneous
subfield embeddings into the faithful four-dimensional Suzuki model.
Two classical facts are retained: $S_d$ is nonabelian simple for
$d\ge3$, and $S_d$ is maximal in $S_e$ if $e/d$ is prime and $d>1$.
These are among the Suzuki structural inputs used by the source itself
(its Table 1, printed page 6).

We also record why subfield intersections are exact in these embeddings.
In characteristic two the standard root and diagonal matrices are
\[
 U(a,b)=\begin{pmatrix}
 1&0&0&0\\a&1&0&0\\a^{\theta+1}+b&a^\theta&1&0\\
 a^{\theta+2}+ab+b^\theta&b&a&1
 \end{pmatrix},\qquad
 h(t)=\operatorname{diag}(t^{\theta+1},t,t^{-1},t^{-\theta-1}),
\]
where $\theta=2^{(e+1)/2}$ in $\F_{2^e}$; let $w$ be the
antidiagonal identity matrix. The standard Bruhat forms are
$U(a,b)h(t)$ and $U(a,b)h(t)wU(c,f)$. This matrix description is
retained classical Suzuki theory. For odd $d\mid e$, the exponent
$\theta$ restricts to $2^{(d+1)/2}$ on $\F_{2^d}$.

\begin{lemma}\label{intersection}
For odd divisors $a,b$ of an odd integer $n$, the standard embedded
subgroups satisfy
\[
                       S_a\cap S_b=S_{\gcd(a,b)}.
\]
\end{lemma}
\begin{proof}
First, the elements of $S_n$ whose matrix entries lie in $\F_{2^d}$
are exactly $S_d$. In the first Bruhat form, the entry $g_{22}$ gives
$t$, the ratio $g_{21}/g_{11}$ gives $a$, and $g_{42}/t$ gives $b$.
Thus all parameters lie in the smaller field when the entries do.
In the second form,
\[
 g_{14}=t^{\theta+1},\quad c=g_{13}/g_{14},\quad
 f=g_{12}/g_{14},\quad a=g_{24}/g_{14},\quad
 b=g_{34}/g_{14}+a^{\theta+1}.
\]
The power $\theta+1$ is coprime to $2^n-1$: an odd common divisor
would give both $2^{n+1}=1$ and $2^n=1$ modulo that divisor.
Consequently $t$ is a fixed integer power of $g_{14}$, and hence also
lies in the smaller field. The two forms are distinguished by whether
$g_{14}$ is zero. This proves the asserted equality with the matrices
over $\F_{2^d}$. Finally
$\F_{2^a}\cap\F_{2^b}=\F_{2^{\gcd(a,b)}}$, which gives the result.
\end{proof}

\section{The omitted pregeometry}
Let $p_1,\ldots,p_k$ be distinct prime divisors of an odd integer $n$,
and put
\[
 d=\frac{n}{p_1\cdots p_k},\qquad G=S_n,\qquad G_i=S_{n/p_i}.
\]
All embeddings are the standard ones just fixed.

\begin{theorem}\label{family}
If $d>1$, the resulting rank-$k$ coset pregeometry is residually
weakly primitive. In fact every available one-step inclusion in its
intersection lattice is maximal, and all its $2^k$ intersections are
distinct.
\end{theorem}
\begin{proof}
Repeated use of Lemma~\ref{intersection} gives, for every $J\subseteq I$,
\[
                G_J=S_{n_J},\qquad n_J=\frac{n}{\prod_{j\in J}p_j}.
\]
For $i\notin J$ the ratio $n_J/n_{J\cup\{i\}}$ is the prime $p_i$.
The smaller exponent is at least $d>1$. Classical subfield maximality
therefore applies to every such inclusion. Unique factorization makes
the integers $n_J$ distinct for distinct subsets $J$. The order
$2^{2r}(2^{2r}+1)(2^r-1)$ of $S_r$ is strictly increasing with $r$,
so the corresponding subgroups are distinct as well. These statements
give every condition in the source's definition.
\end{proof}

For an explicit example take
\[
       n=3^2\cdot5\cdot7\cdot11=3465,
       \qquad (p_1,p_2,p_3,p_4)=(3,5,7,11).
\]
The four maximal parabolics are
\[
                 S_{1155},\quad S_{693},\quad S_{495},\quad S_{315},
\]
and their common intersection is $S_3=\Sz(8)$.
Theorem~\ref{family} proves that this is a rank-four residually weakly
primitive pregeometry.

\begin{theorem}\label{omission}
The displayed pregeometry is absent from Constructions 1--11 in the
fixed published edition of \emph{The residually weakly primitive
pre-geometries of the Suzuki simple groups}. It therefore contradicts
the literal exhaustive assertion in its Theorem 5.
\end{theorem}
\begin{proof}
Constructions 1--10 all include a soluble maximal parabolic: a dihedral
group, an affine soluble group, a cyclic group extended by a group of
order four, or the specified mixed root group. Every maximal parabolic
in our example is nonabelian simple. Thus a permutation of types and
an ambient group automorphism cannot identify our system with any of
those constructions.

Construction 11 writes $n=p_1^{e_1}\cdots p_m^{e_m}$ with the $p_i$
pairwise distinct, and permits ranks $\alpha\in\{2,\ldots,m-1\}$.
Here $m=4$, so its largest permitted rank is three. Our rank is four.
Theorem 5 states that all residually weakly primitive pregeometries of
rank at least four are those in Constructions 1--11, with no additional
operation of taking truncations. The example is therefore outside that
literal list.
\end{proof}

\section{Attribution and limits}
The source is D.~Leemans, \emph{The residually weakly primitive
pre-geometries of the Suzuki simple groups}, Note di Matematica 20(1)
(2000/2001), 1--20. The relevant definition is on printed pages 4--5,
the construction conventions on page 15, and Construction 11 and
Theorem 5 on page 18. The source was received on 26 July 1999 and
accepted on 21 September 1999.

The argument is a reworking of Sections 2--3 of the project's earlier
Suzuki subfield manuscript, whose separate review already checked this
exact omitted boundary against the source. The old proof establishes
additional geometry and flag-transitivity statements using descent.
Those stronger conclusions are not needed here and are not acquired
as premises of this note. The elementary matrix intersection argument
comes from the earlier project treatment of compatible Suzuki tori.
Classical Suzuki structure and prime-degree subfield maximality remain
explicit retained inputs.

The correction concerns the literal rank range. The condition
\[
                         n/(p_1\cdots p_k)>1
\]
includes the example; suitable extensions
of the construction can therefore repair this omission. We make no
claim that such a repair exhausts all embeddings, and do not classify
the remaining ten constructions. Neither a new mathematical discovery
nor historical priority is claimed by this reworking. It contributes
no additional conjecture solution, Lean theorem or compression factor.
The supplementary control checks the matrix parameter identities
symbolically, and the sixteen exponents and thirty-two prime-degree
inclusions by integer arithmetic. It does not enumerate $S_{3465}$.
\end{document}
